% Redrawn 2026-10-15; main labels 8 pt, notes 7.5 pt; actual fonts embedded in PDF.
% export-reference-width-mm: 169
\begin{tikzpicture}[font=\figurefont\fontsize{8}{10}\selectfont,
  box/.style={line width=.6pt,rounded corners=1.8pt,draw=FigureStroke,fill=FigureBlueFill,align=center,
    text width=30mm,minimum height=14mm},
  arr/.style={-{Latex[length=1.8mm]},draw=FigureStroke,line width=.6pt},
  control/.style={draw=BookAmber,dashed,-{Latex[length=1.8mm]}}]
  \node[box,fill=FigureGreenFill] (model) at (0,0) {学习模型\\$\widehat T,\widehat R$};
  \node[box,fill=FigureYellowFill] (policy) at (5,0) {策略改进\\改变访问分布};
  \node[box,fill=FigureBlueFill] (env) at (10,0) {真实执行与校验\\获得新交互记录};
  \node[box,fill=FigureGreenFill] (fit) at (5,-2.8) {在新访问区域\\重新拟合模型};
  \draw[arr] (model) -- (policy);
  \draw[arr] (policy) -- (env);
  \draw[arr] (env.south) |- (fit.east);
  \draw[arr] (fit.west) -| (model.south);
  \node[align=center,text width=29mm] (penalty) at (0,2.3) {不确定性惩罚\\改变模型评分};
  \node[align=center,text width=29mm] (short) at (5,2.3) {缩短想象时域\\限制传播深度};
  \node[align=center,text width=29mm] (cal) at (10,2.3) {新策略交互记录校准\\需要额外交互};
  \draw[control] (penalty) -- (model);
  \draw[control] (short) -- (policy);
  \draw[control] (cal) -- (env);
  \node[align=center,text width=128mm] at (5,-4.3)
    {离线条件下右侧的真实交互边不可凭空补上。\\固定策略上的低误差，不能自动控制新策略的误差。};
\end{tikzpicture}
